% Edit this file first. Keep the YYYY-MM-DD date form used by the standard.
\ReportNumber{TR-2026-001}
\ReportSecurity{公开}
\ReportSecurityTerm{}
\ReportIssueDate{2026-08-21}

\ReportType{专题报告}
\ReportTitle{示例专题科技报告：复杂技术系统的验证方法}
\ReportEnglishTitle{A Verification Method for Complex Technical Systems}
\ReportVolume{}
\ReportEnglishVolume{}

\ReportAuthor{张三，李四}
\ReportInstitution{示例研究院}
\ReportCompletionDate{2026-08-21}
\ReportPeriod{2026-01-01--2026-08-21}

\ReportNote{本文件为 GB/T 7713.3-2014 模板的排版示例，内容仅用于演示。}

\ReportAbstractCN{本报告面向复杂技术系统的专题研究与工程验证场景，提出一套可复核、可追踪的验证方法及科技报告组织方案。首先，依据任务目标识别系统边界、关键假设、性能指标和风险约束，形成从研究问题到证据材料的双向追踪关系；其次，采用理论分析、受控试验、数据核验和不确定性分析相结合的方法，对模型、算法、装置或软件系统的有效性进行分层验证；再次，以统一编号的图、表、公式和附录记录输入条件、试验过程、异常结果及处置依据，并通过独立复算检查主要结论。示例结果表明，该方法能够减少验证口径不一致、证据链断裂和结论不可复现等问题，使专业读者在掌握必要条件的情况下复核研究过程。报告同时给出实施建议，包括在项目早期确定元数据责任人、统一计量单位和术语、保留正反两方面结果，以及在提交前执行内容、格式和引用三类检查。该方案可作为专题科技报告撰写与质量控制的起点，但具体指标和专业规范仍应由项目主管部门结合所属领域进一步确定。}
\ReportKeywordsCN{科技报告；专题报告；系统验证；证据追踪；质量控制}

\ReportAbstractEN{This report presents a reproducible and traceable verification method for topic-oriented research on complex technical systems, together with a document structure suitable for a scientific and technical report. The method begins by defining the system boundary, research questions, assumptions, performance indicators, decision criteria, and risk constraints. Each requirement is then linked to evidence that can be inspected independently, including theoretical derivations, controlled experiments, source data, calculation records, software versions, and review notes. Verification is organized in layers so that models, algorithms, devices, and integrated systems can be assessed at an appropriate level before conclusions are combined. The procedure uses analytical checks, repeatable experiments, data validation, sensitivity analysis, and uncertainty analysis. Inputs, operating conditions, abnormal observations, negative results, and corrective actions are retained rather than removed from the record. Figures, tables, equations, and appendices receive consistent identifiers so that readers can move from a claim to the supporting material and back again. In the illustrative application, the method reduces inconsistent evaluation criteria, broken evidence chains, and results that cannot be reproduced. It also improves the separation between observations, interpretations, conclusions, and recommendations. Practical recommendations include assigning responsibility for report metadata at the beginning of the project, using statutory measurement units and controlled terminology, checking every citation against the reference list, and conducting separate reviews of technical content, document structure, and typesetting before submission. The proposed workflow is intended as a starting point rather than a substitute for discipline-specific rules. Project sponsors and competent authorities should define additional indicators, confidentiality markings, data-retention requirements, and acceptance procedures for their own fields. When those local requirements conflict with optional presentation choices in this template, the approved project or institutional rules should take precedence while the mandatory structural elements of the national standard are retained.}
\ReportKeywordsEN{scientific and technical report; topic report; system verification; evidence traceability; quality control}

\ReportSupportChannel{国家重点研发计划（示例）}
\ReportProjectName{复杂技术系统验证与评价方法研究}
\ReportContractor{示例研究院}
\ReportProjectLeader{王五}
\ReportSponsor{示例主管部门}
\ReportProjectNumber{2026YF000001}
\ReportProgramName{技术创新专项}
\ReportContactName{赵六}
\ReportContactDetails{010-00000000}

\ReportPublisherPlace{北京}
\ReportPublisher{示例研究院}
\ReportPublisherDate{2026-08-21}
